% Unofficial New York University Poster Template
% a fork of https://github.com/anishathalye/gemini
% also refer to https://github.com/k4rtik/uchicago-poster

% Adapted by Phillip Wang, all credits go to Tao Li


\documentclass[final]{beamer}

% ====================
% Packages
% ====================

\usepackage[T1]{fontenc}
%\usepackage[utf8]{luainputenc}
\usepackage{lmodern}
\usepackage[size=custom, width=59.4,height=42, scale=0.7]{beamerposter}
\usetheme{gemini}
\usecolortheme{cam}
\usepackage{graphicx}
\usepackage{booktabs}
\usepackage{tikz}
\usepackage{pgfplots}
\pgfplotsset{compat=1.14}
\usepackage{anyfontsize}

% ====================
% Lengths
% ====================

% If you have N columns, choose \sepwidth and \colwidth such that
% (N+1)*\sepwidth + N*\colwidth = \paperwidth
\newlength{\sepwidth}
\newlength{\colwidth}
\setlength{\sepwidth}{0.025\paperwidth}
\setlength{\colwidth}{0.3\paperwidth}

% Block without a title, used to continue a block from the previous column
\newenvironment{contblock}
  {\begin{beamercolorbox}[colsep*=0ex]{block body}%
    \usebeamerfont{block body}\justifying\setlength{\parskip}{1ex}}
  {\end{beamercolorbox}\vskip0pt\vspace*{2ex}}

\newcommand{\separatorcolumn}{\begin{column}{\sepwidth}\end{column}}

% ====================
% Title
% ====================

\title{Title}

\author{Department of Physics, Durham University}

\institute[shortinst]{Short abstract.}

% ====================
% Footer (optional)
% ====================

\footercontent{
  %\href{https://github.com}{Github: https://github.com} \hfill
  %Capstone Festival, Abu Dhabi \hfill
  %\href{mailto:taoli@nyu.edu}{faiza@nyu.edu}}
}
% (can be left out to remove footer)

% ====================
% Logo (optional)
% ====================

% use this to include logos on the left and/or right side of the header:
% Left: institution
 \logoleft{\includegraphics[height=5cm]{logos/durham-university-logo-transparent.png}}
% Right: funding agencies and other affilations
%\logoright{\includegraphics[height=7cm]{logos/NSF.eps}}
% ====================
% Body
% ====================

\begin{document}



\begin{frame}[t]
\begin{columns}[t]
\separatorcolumn

\begin{column}{\colwidth}

  \begin{block}{Atomic Qubits with Rydberg States}

    The young field of quantum computing is expected by many to revolutionise various aspects of society. These include
    cryptography, drug and materials research, optimisation and more. It works by engangling quantum bits (qubits)
    in useful ways. There are various choices for the implementation of qubits in practice, one of which being an atom with a
    low-excited state
    and a high-excited state coupled by an electromagnetic (EM) field. This investigation probes the dynamics of a quantum system
    composed of two-such qubits.

    Consider an atomic qubit driven by a plane-polarised EM wave that perturbs the atom's original hamiltonian $H_0$ by
     $H'(t)=-\mathbf{d} \cdot \mathbf{\epsilon}$,
    where $\mathbf{d}$
    is the dipole operator and $\mathbf{\epsilon}$ is a sinusoidally-time-varying electric field. The qubit's time evolution is
    governed by a time-dependent Schr\"odinger equation, $i\hbar \partial_t |\psi(t)\rangle = H_{int} |\psi(t)\rangle$, where
    $|\psi(t)\rangle$ is the qubit's state vector and $H_{int}$ describes the interaction between the atom and the EM field. Let
    \{$|0\rangle, |r\rangle$\} be a set of orthonormal eigenstates of $H_0$ such that $|0\rangle$ is a low-excited state with
    energy $E_0$ and
    $|r\rangle$ is a high-excited state with principal quantum number on the order of $n>50$ and energy $E_r$,
    known as a Rydberg state. The
    natural frequency of the transition between these states is $\omega_0 = (E_r - E_0) / \hbar$, and the Rabi frequency is
    defined as $\Omega=\langle r|H'(t)|0\rangle / \hbar$.  When
    considering interaction with an EM wave of angular frequency $\omega$ and detuning $\Delta = \omega - \omega_0$
    such that $|\Delta|,\Omega \ll |\omega - \omega_{other}|$, where $\omega_{other}$ ranges over the frequencies of all other
    laser-accessible transitions, $|\psi\rangle$ can be approximated as a superposition of only the eigenstates
    \{$|0\rangle, |r\rangle$\}. In this basis the atomic state vector and interaction hamiltonian are
    $$
    |\psi(t)\rangle = \begin{pmatrix}
    a(t) \\
    b(t) \\
    \end{pmatrix}, H_{int} = \hbar/2
    \begin{pmatrix}
    \Delta & \Omega e^{-i \phi_L} \\
    \Omega e^{i \phi_L} & -\Delta
    \end{pmatrix},
    $$
    where $\phi_L$ is the initial phase of the electric field, and $a(t),b(t)$ are time-dependent complex coefficients.

    Any useful quantum computation requires entanglement between multiple qubits, however neutrally charged atoms do not
    produce a strong enough interaction to be entangled; this is why the atomic qubit described involves a Rydberg state. When
    two nearby atoms are in Rydberg states, their van der Waals interaction causes the combined state to experience
    an energy shift of $\hbar V = C_6 / R^6$, where $R$ is their internuclear distance and $C_6$ is a constant rising rapidly
    with principle quantum number. In the two-atom basis \{$|00\rangle, |0r\rangle, |r0\rangle, |rr\rangle$\}, representing
    the four eigenstates of the bare two-atom hamiltonian, the two-atom state vector and interaction hamiltonian are






  \end{block}





\end{column}

\separatorcolumn

\begin{column}{\colwidth}

  $$ |\psi(t)\rangle = \begin{pmatrix}
    a(t) \\
    b(t) \\
    c(t) \\
    d(t)
    \end{pmatrix}, H_{int}= \hbar / 2
    \begin{pmatrix}
    2\Delta & \Omega e^{-i \phi_L} & \Omega e^{-i \phi_L} & 0 \\
    \Omega e^{i \phi_L} & 0 & 0 & \Omega e^{-i \phi_L} \\
    \Omega e^{i \phi_L} & 0 & 0 & \Omega e^{-i \phi_L} \\
    0 & \Omega e^{i \phi_L} & \Omega e^{i \phi_L} & -2\Delta + 2V
    \end{pmatrix}.
    $$

  \begin{block}{A block containing an enumerated list}


    Vivamus congue volutpat elit non semper. Praesent molestie nec erat ac
    interdum. In quis suscipit erat. \textbf{Phasellus mauris felis, molestie
    ac pharetra quis}, tempus nec ante. Donec finibus ante vel purus mollis
    fermentum. Sed felis mi, pharetra eget nibh a, feugiat eleifend dolor. Nam
    mollis condimentum purus quis sodales. Nullam eu felis eu nulla eleifend
    bibendum nec eu lorem. Vivamus felis velit, volutpat ut facilisis ac,
    commodo in metus.

    \begin{enumerate}
      \item \textbf{Morbi mauris purus}, egestas at vehicula et, convallis
        accumsan orci. Orci varius natoque penatibus et magnis dis parturient
        montes, nascetur ridiculus mus.
      \item \textbf{Cras vehicula blandit urna ut maximus}. Aliquam blandit nec
        massa ac sollicitudin. Curabitur cursus, metus nec imperdiet bibendum,
        velit lectus faucibus dolor, quis gravida metus mauris gravida turpis.
      \item \textbf{Vestibulum et massa diam}. Phasellus fermentum augue non
        nulla accumsan, non rhoncus lectus condimentum.
    \end{enumerate}

  \end{block}

  \begin{block}{Fusce aliquam magna velit}

    Et rutrum ex euismod vel. Pellentesque ultricies, velit in fermentum
    vestibulum, lectus nisi pretium nibh, sit amet aliquam lectus augue vel
    velit. Suspendisse rhoncus massa porttitor augue feugiat molestie. Sed
    molestie ut orci nec malesuada. Sed ultricies feugiat est fringilla
    posuere.

    \begin{figure}
      \centering
      \begin{tikzpicture}
        \begin{axis}[
            scale only axis,
            no markers,
            domain=0:2*pi,
            samples=100,
            axis lines=center,
            axis line style={-},
            ticks=none]
          \addplot[red] {sin(deg(x))};
          \addplot[blue] {cos(deg(x))};
        \end{axis}
      \end{tikzpicture}
      \caption{Another figure caption.}
    \end{figure}

  \end{block}

  \begin{block}{Nam cursus consequat egestas}

    Nulla eget sem quam. Ut aliquam volutpat nisi vestibulum convallis. Nunc a
    lectus et eros facilisis hendrerit eu non urna. Interdum et malesuada fames
    ac ante \textit{ipsum primis} in faucibus. Etiam sit amet velit eget sem
    euismod tristique. Praesent enim erat, porta vel mattis sed, pharetra sed
    ipsum. Morbi commodo condimentum massa, \textit{tempus venenatis} massa
    hendrerit quis. Maecenas sed porta est. Praesent mollis interdum lectus,
    sit amet sollicitudin risus tincidunt non.

    Etiam sit amet tempus lorem, aliquet condimentum velit. Donec et nibh
    consequat, sagittis ex eget, dictum orci. Etiam quis semper ante. Ut eu
    mauris purus. Proin nec consectetur ligula. Mauris pretium molestie
    ullamcorper. Integer nisi neque, aliquet et odio non, sagittis porta justo.

    \begin{itemize}
      \item \textbf{Sed consequat} id ante vel efficitur. Praesent congue massa
        sed est scelerisque, elementum mollis augue iaculis.
        \begin{itemize}
          \item In sed est finibus, vulputate
            nunc gravida, pulvinar lorem. In maximus nunc dolor, sed auctor eros
            porttitor quis.
          \item Fusce ornare dignissim nisi. Nam sit amet risus vel lacus
            tempor tincidunt eu a arcu.
          \item Donec rhoncus vestibulum erat, quis aliquam leo
            gravida egestas.
        \end{itemize}
      \item \textbf{Sed luctus, elit sit amet} dictum maximus, diam dolor
        faucibus purus, sed lobortis justo erat id turpis.
      \item \textbf{Pellentesque facilisis dolor in leo} bibendum congue.
        Maecenas congue finibus justo, vitae eleifend urna facilisis at.
    \end{itemize}

  \end{block}

\end{column}

\separatorcolumn

\begin{column}{\colwidth}

  \begin{exampleblock}{A highlighted block containing some math}

    A different kind of highlighted block.

    $$
    \int_{-\infty}^{\infty} e^{-x^2}\,dx = \sqrt{\pi}
    $$

    Interdum et malesuada fames $\{1, 4, 9, \ldots\}$ ac ante ipsum primis in
    faucibus. Cras eleifend dolor eu nulla suscipit suscipit. Sed lobortis non
    felis id vulputate.

    \heading{A heading inside a block}

    Praesent consectetur mi $x^2 + y^2$ metus, nec vestibulum justo viverra
    nec. Proin eget nulla pretium, egestas magna aliquam, mollis neque. Vivamus
    dictum $\mathbf{u}^\intercal\mathbf{v}$ sagittis odio, vel porta erat
    congue sed. Maecenas ut dolor quis arcu auctor porttitor.

    \heading{Another heading inside a block}

    Sed augue erat, scelerisque a purus ultricies, placerat porttitor neque.
    Donec $P(y \mid x)$ fermentum consectetur $\nabla_x P(y \mid x)$ sapien
    sagittis egestas. Duis eget leo euismod nunc viverra imperdiet nec id
    justo.

  \end{exampleblock}

  \begin{block}{References}

    Class aptent taciti sociosqu ad litora torquent per conubia nostra, per
    inceptos himenaeos. Phasellus libero enim, gravida sed erat sit amet,
    scelerisque congue diam. Fusce dapibus dui ut augue pulvinar iaculis.

    \begin{table}
      \centering
      \begin{tabular}{l r r c}
        \toprule
        \textbf{First column} & \textbf{Second column} & \textbf{Third column} & \textbf{Fourth} \\
        \midrule
        Foo & 13.37 & 384,394 & $\alpha$ \\
        Bar & 2.17 & 1,392 & $\beta$ \\
        Baz & 3.14 & 83,742 & $\delta$ \\
        Qux & 7.59 & 974 & $\gamma$ \\
        \bottomrule
      \end{tabular}
      \caption{A table caption.}
    \end{table}

    Donec quis posuere ligula. Nunc feugiat elit a mi malesuada consequat. Sed
    imperdiet augue ac nibh aliquet tristique. Aenean eu tortor vulputate,
    eleifend lorem in, dictum urna. Proin auctor ante in augue tincidunt
    tempor. Proin pellentesque vulputate odio, ac gravida nulla posuere
    efficitur. Aenean at velit vel dolor blandit molestie. Mauris laoreet
    commodo quam, non luctus nibh ullamcorper in. Class aptent taciti sociosqu
    ad litora torquent per conubia nostra, per inceptos himenaeos.



  \end{block}

  \begin{block}{References}

    \nocite{*}
    \footnotesize{\bibliographystyle{plain}\bibliography{poster}}

  \end{block}

\end{column}

\separatorcolumn
\end{columns}
\end{frame}

\end{document}
